% =====================================================================
%  Основное содержимое отчёта.
% =====================================================================

\labpart{Описание программы}

Программа реализована на языке C\# (.NET 9, консольное приложение) и
последовательно выполняет все пункты задания: создание учётных записей и
групп, создание каталога \texttt{pzs} с пятью вложенными каталогами под
разные классы субъектов доступа, создание файловой матрицы прав в каждом
из них, проверку прав доступа для каждой пары «пользователь -- объект»,
запуск и попытку остановки процессов, проверку операций над каталогами и
финальную очистку. Поскольку программа создаёт настоящих пользователей и
групп и меняет права доступа реальной файловой системы, запускать её
можно только от имени \texttt{root} -- это проверяется в самом начале
работы. Между шагами программа приостанавливается и ждёт нажатия клавиши,
что позволяет в любой момент открыть отдельный терминал и убедиться в
реальном состоянии системы (\texttt{ls~-la}, \texttt{getfacl}, \texttt{id}
и т.\,п.), прежде чем перейти к следующему шагу.

Классические биты \texttt{chmod} достаточны, чтобы выдать права владельцу,
«остальным» или всем сразу, однако они не позволяют выдать права ровно
одной именованной группе, не затронув при этом остальных, -- группа файла
в POSIX всего одна. Поэтому для класса «группа» классические биты
обнуляются (\texttt{chmod~000}), а затем командой \texttt{setfacl}
добавляется именованная ACE для \texttt{group\_iit1}: это и есть
собственно список контроля доступа, а не просто традиционные три класса
прав. Соответствующий фрагмент кода показан ниже (см. \autoref{lst:rules}).

\begin{codelisting}[fontsize=\footnotesize]{csharp}{Применение правила доступа к объекту файловой системы}{rules}
public static void ApplyRule(string path, int subject, int permIndex)
{
    var octal = Config.PermOctals[permIndex];
    var sym = Config.PermSyms[permIndex];

    Report.Execute($"Reset ACL on '{path}'", "setfacl", "-b", path);

    switch (subject)
    {
        case 1: // owner
            Report.Execute($"Apply access rule (owner, {sym}) on '{path}'", "chmod", $"{octal}00", path);
            break;

        case 2: // named group group_iit1
            Report.Execute($"Clear classic bits on '{path}'", "chmod", "000", path);
            Report.Execute($"Apply access rule (group 'group_iit1', {sym}) on '{path}'",
                "setfacl", "-m", $"g:group_iit1:{sym}", path);
            break;

        case 3: // other
            Report.Execute($"Apply access rule (other, {sym}) on '{path}'", "chmod", $"00{octal}", path);
            break;

        case 4: // all
            Report.Execute($"Apply access rule (all, {sym}) on '{path}'", "chmod", $"{octal}{octal}{octal}", path);
            break;

        case 5: // admin (root) only
            Report.Execute($"Set owner of '{path}' to root", "chown", "root:root", path);
            Report.Execute($"Apply access rule (administrator only, {sym}) on '{path}'", "chmod", $"{octal}00", path);
            break;
    }
}
\end{codelisting}

\labpartbreak{Результаты выполнения}

Ниже приведён фрагмент реального вывода программы, запущенной в
изолированном Linux-контейнере (см. \autoref{lst:output}). Он
показывает, что каталог \texttt{pzs12} доступен на чтение/запись только
членам группы \texttt{group\_iit1} (\texttt{iit11}, \texttt{iit12}), тогда
как \texttt{iit21}, \texttt{iit22}, \texttt{iit3} получают отказ, а
\texttt{root} проходит любую проверку.

Отдельного внимания заслуживает файл \texttt{pzs14/file21}: он принадлежит
\texttt{iit11}, а право на чтение по ACL выдано группе
\texttt{group\_iit1}, в которую \texttt{iit11} тоже входит. Тем не менее
\texttt{iit11} получает отказ, а \texttt{iit12} -- доступ. Это в точности
соответствует правилу выбора ACE, описанному в теоретической части
задания: если запрашивающий -- владелец объекта, то используется только
запись владельца (обнулённая \texttt{chmod~000}), а именованная запись
группы для него не рассматривается вовсе, даже если он формально состоит
в этой группе.

Также подтвердилось, что ни один из файлов вида \texttt{filex5}
(право только на выполнение, без права на чтение) не удаётся запустить ни
одному пользователю, включая владельца \texttt{iit11}: интерпретируемый
sh/bash-скрипт запускается не самим ядром, а внешним интерпретатором,
которому требуется прочитать содержимое файла, -- права на одно только
выполнение для этого недостаточно (в отличие от скомпилированных
бинарных файлов).

\codelistingfile[fontsize=\footnotesize]{text}{lab/code/output.txt}{Фрагмент вывода программы}{output}
